\documentclass[11pt,aspectratio=169]{beamer}
\usetheme{SimplePlus}

\usepackage[T1]{fontenc}
\usepackage[utf8]{inputenc}
\usepackage{amsmath,amssymb,mathtools}
\usepackage{booktabs}
\usepackage{array}
\usepackage{appendixnumberbeamer}
\usepackage{hyperref}
\pdfstringdefDisableCommands{\def\translate#1{#1}}

\newcommand{\E}{\mathbb{E}}
\newcommand{\Prob}{\mathbb{P}}
\newcommand{\Var}{\mathrm{Var}}
\newcommand{\Cov}{\mathrm{Cov}}
\newcommand{\ATT}{\mathrm{ATT}}
\newcommand{\ATE}{\mathrm{ATE}}
\newcommand{\ATU}{\mathrm{ATU}}
\newcommand{\argmin}{\mathop{\mathrm{arg\,min}}}
\newcommand{\indep}{\perp\!\!\!\perp}
\newcommand{\indicator}{\mathbf{1}}
\newcommand{\pscore}{e(X_i)}
\newcommand{\derivbutton}[1]{\vspace{0.4em}\hfill\hyperlink{app:#1}{\beamerbutton{Appendix details}}}
\newcommand{\backbutton}[1]{\vspace{0.4em}\hfill\hyperlink{main:#1}{\beamerbutton{Back to main slide}}}

\title{Applied Econometrics: Session 6}
\subtitle{Matching and Propensity Scores}
\author{PhD Intensive Course}
\institute{Lecture 6}
\date{\today}

\begin{document}

\begin{frame}
  \titlepage
\end{frame}

\begin{frame}{Why This Lecture Matters}
  Many empirical papers have a treatment group and a comparison group, but no random assignment.
  \vspace{0.4em}

  Matching and propensity scores ask:
  \begin{block}{Core question}
    Can we make treated and untreated units comparable by conditioning carefully on pre-treatment observables?
  \end{block}

  \begin{itemize}
    \item Matching tries to compare like with like.
    \item Propensity scores compress many observed covariates into one treatment probability.
    \item Neither method fixes selection on unobservables.
  \end{itemize}
\end{frame}

\begin{frame}{Roadmap for Today}
  \tableofcontents
\end{frame}

\section{Setup}

\begin{frame}{The Observational Problem}
  \hypertarget{main:setup}{}
  We observe
  \[
    (Y_i,D_i,X_i), \qquad D_i\in\{0,1\}.
  \]
  \begin{itemize}
    \item \(Y_i\): observed outcome.
    \item \(D_i\): treatment indicator.
    \item \(X_i\): pre-treatment covariates, such as age, education, past outcomes, location.
  \end{itemize}
  Potential outcomes:
  \[
    Y_i = D_iY_{1i}+(1-D_i)Y_{0i}.
  \]
  The problem is that treated and untreated people may differ in \(X_i\) before treatment.
  \derivbutton{cia-matching}
\end{frame}

\begin{frame}{The Required Assumption}
  \hypertarget{main:cia}{}
  Matching and propensity-score methods usually rely on conditional independence:
  \[
    (Y_{1i},Y_{0i})\indep D_i \mid X_i.
  \]
  \begin{itemize}
    \item After comparing units with the same \(X_i\), treatment is as good as random.
    \item This is a strong economic assumption, not a theorem.
    \item \(X_i\) must be measured before treatment.
    \item Post-treatment variables are usually bad controls.
  \end{itemize}
  \begin{block}{Plain English}
    Among people with the same observed background, the treated and untreated differ only because of treatment.
  \end{block}
  \derivbutton{cia-matching}
\end{frame}

\begin{frame}{Overlap: The Other Required Assumption}
  \hypertarget{main:overlap}{}
  Conditional independence is not enough. We also need common support:
  \[
    0 < \Prob(D_i=1\mid X_i=x) < 1
    \qquad \text{for relevant } x.
  \]
  \begin{itemize}
    \item For each treated type, we need comparable untreated units.
    \item For each untreated type, we need comparable treated units if the target is the ATE.
    \item No estimator can learn a counterfactual from a covariate region containing only one treatment status.
  \end{itemize}
  \begin{block}{Practical warning}
    Bad overlap is not a software problem. It changes what population can be studied credibly.
  \end{block}
  \derivbutton{overlap-trimming}
\end{frame}

\begin{frame}{Which Effect Are We Estimating?}
  Three common targets:
  \[
    \ATE=\E[Y_{1i}-Y_{0i}],
  \]
  \[
    \ATT=\E[Y_{1i}-Y_{0i}\mid D_i=1],
  \]
  \[
    \ATU=\E[Y_{1i}-Y_{0i}\mid D_i=0].
  \]
  \begin{itemize}
    \item Matching treated units to controls naturally estimates the ATT.
    \item Symmetric matching or IPW can target the ATE.
    \item With heterogeneous treatment effects, these are not the same object.
  \end{itemize}
\end{frame}

\section{Matching}

\begin{frame}{Matching: The Main Idea}
  \begin{itemize}
    \item For each treated unit, find untreated units that look similar before treatment.
    \item Use those untreated outcomes as estimates of the treated unit's missing \(Y_{0i}\).
    \item Then compare observed treated outcomes to imputed untreated outcomes.
  \end{itemize}
  \begin{block}{Matching as design}
    We are trying to construct a comparison sample before looking too hard at outcomes.
  \end{block}
  \begin{block}{Matching as estimation}
    We are imputing the missing potential outcome using observed units with similar \(X_i\).
  \end{block}
\end{frame}

\begin{frame}{Exact Matching}
  Exact matching compares treated and untreated units in the same covariate cell.
  \[
    \delta(x)
    =
    \E[Y_i\mid D_i=1,X_i=x]
    -
    \E[Y_i\mid D_i=0,X_i=x].
  \]
  The ATT is a weighted average over treated cells:
  \[
    \ATT
    =
    \sum_x \delta(x)\Prob(X_i=x\mid D_i=1).
  \]
  \begin{itemize}
    \item Works well when \(X_i\) is low-dimensional and discrete.
    \item Fails quickly when \(X_i\) has many categories or continuous variables.
  \end{itemize}
  \derivbutton{exact-matching}
\end{frame}

\begin{frame}{A Tiny Exact-Matching Example}
  \scriptsize
  \begin{tabular}{>{\raggedright\arraybackslash}p{0.16\linewidth}ccccc}
    \toprule
    Cell \(X\) & Treated & Control & Effect & Treated share & Contribution \\
    \midrule
    Young, urban & 12.0 & 9.5 & 2.5 & 0.40 & 1.00 \\
    Young, rural & 10.0 & 8.0 & 2.0 & 0.20 & 0.40 \\
    Older, urban & 14.0 & 12.5 & 1.5 & 0.25 & 0.38 \\
    Older, rural & 11.0 & 10.0 & 1.0 & 0.15 & 0.15 \\
    \midrule
    \multicolumn{5}{r}{Exact-matching ATT} & 1.93 \\
    \bottomrule
  \end{tabular}
  \vspace{0.6em}

  \begin{itemize}
    \item We do not compare all treated to all controls.
    \item We compare within covariate cells and then average using the treated distribution.
    \item The weights matter because the target is the treated population.
  \end{itemize}
\end{frame}

\begin{frame}{Matching as Imputation}
  \hypertarget{main:imputation}{}
  For a treated unit \(i\), \(Y_{1i}\) is observed but \(Y_{0i}\) is missing.
  Matching estimates
  \[
    \widehat{Y}_{0i}
    =
    \frac{1}{K}\sum_{\ell\in\mathcal{N}_K(i;D=0)}Y_\ell,
  \]
  where \(\mathcal{N}_K(i;D=0)\) is the set of \(K\) nearest untreated neighbors.
  Then
  \[
    \widehat{\ATT}_{M}
    =
    \frac{1}{N_1}\sum_{i:D_i=1}(Y_i-\widehat{Y}_{0i}).
  \]
  \begin{itemize}
    \item The missing counterfactual is filled in using nearby controls.
    \item Better neighbors reduce bias; more neighbors can reduce variance.
  \end{itemize}
  \derivbutton{nn-imputation}
\end{frame}

\begin{frame}{Nearest-Neighbor Matching}
  Exact matching is often impossible, so we use a distance.
  \[
    d(i,\ell)
    =
    \left[(X_i-X_\ell)'W(X_i-X_\ell)\right]^{1/2}.
  \]
  \begin{itemize}
    \item \(W=I\): Euclidean distance after standardizing covariates.
    \item \(W=\widehat{\Sigma}_X^{-1}\): Mahalanobis distance.
    \item The closest untreated unit becomes the match for a treated unit.
  \end{itemize}
  \begin{block}{Important detail}
    Always scale variables before distance matching. Otherwise income in euros can dominate age in years.
  \end{block}
\end{frame}

\begin{frame}{With Replacement or Without Replacement?}
  \small
  \begin{tabular}{>{\raggedright\arraybackslash}p{0.25\linewidth}>{\raggedright\arraybackslash}p{0.32\linewidth}>{\raggedright\arraybackslash}p{0.32\linewidth}}
    \toprule
    Choice & Advantage & Cost \\
    \midrule
    With replacement & Best control can be reused, usually lowers bias & Fewer unique controls, often higher variance \\
    Without replacement & Uses more control observations & Later treated units may receive poor matches \\
    One nearest neighbor & Very local comparison & Noisy estimate \\
    Many neighbors & Lower variance & More distant matches, more bias \\
    Caliper & Blocks bad matches & Drops units outside tolerance \\
    \bottomrule
  \end{tabular}
  \vspace{0.6em}

  There is no universally correct setting. The choice should be reported and checked.
\end{frame}

\begin{frame}{Bias-Variance Tradeoff in Matching}
  \hypertarget{main:biasvariance}{}
  Matching is nonparametric. Its performance depends on how local the comparison is.
  \vspace{0.4em}

  \begin{itemize}
    \item Very strict matching:
      low bias, but many treated units may have no match and estimates are noisy.
    \item Loose matching:
      many matches and lower variance, but treated and control units may not be comparable.
    \item More covariates:
      controls more confounding, but makes good matches harder to find.
  \end{itemize}
  \begin{block}{Practical rule}
    Matching quality is a design object. Inspect distances, covariate balance, and dropped observations.
  \end{block}
  \derivbutton{matching-bias}
\end{frame}

\begin{frame}{The Curse of Dimensionality}
  Suppose each covariate is split into five bins.
  \[
    \text{Number of cells}=5^p,
  \]
  where \(p\) is the number of covariates.
  \vspace{0.4em}

  \begin{tabular}{ccc}
    \toprule
    Covariates \(p\) & Cells \(5^p\) & What happens \\
    \midrule
    2 & 25 & feasible \\
    5 & 3,125 & many empty cells \\
    10 & 9,765,625 & exact matching collapses \\
    \bottomrule
  \end{tabular}
  \vspace{0.5em}

  \begin{block}{Lesson}
    The more dimensions we match on, the farther the nearest neighbor usually is.
  \end{block}
\end{frame}

\begin{frame}{Common Support in Matching}
  \small
  \begin{itemize}
    \item If treated units are older, richer, and more educated than all controls, matching has no credible comparison.
    \item If some treated units have no nearby controls, the ATT for all treated units is not supported by the data.
    \item It may be better to report an effect for the matched or overlapping population.
  \end{itemize}
  \begin{block}{Diagnostic questions}
    \begin{itemize}
      \item Are matched control covariates close to treated covariates?
      \item How many treated units are unmatched or outside the caliper?
      \item Are results driven by a small set of reused controls?
    \end{itemize}
  \end{block}
\end{frame}

\begin{frame}{What Matching Does Not Solve}
  Matching is not magic.
  \begin{itemize}
    \item It does not remove bias from unobserved ability, motivation, risk, or expectations.
    \item It does not make post-treatment controls safe.
    \item It does not prove that all relevant confounders were measured.
    \item It can increase bias if the distance metric prioritizes irrelevant variables.
  \end{itemize}
  \begin{block}{Interpretation}
    A matching estimate is causal only if the remaining assignment variation within matched sets is plausibly as good as random.
  \end{block}
\end{frame}

\section{Propensity Scores}

\begin{frame}{Why Propensity Scores?}
  Matching on many covariates is hard. Propensity scores compress treatment selection into one number:
  \[
    e(X_i)=\Prob(D_i=1\mid X_i).
  \]
  \begin{itemize}
    \item \(e(X_i)\) is the probability of being treated for a unit with covariates \(X_i\).
    \item Units with the same propensity score have the same predicted treatment probability.
    \item We can match, stratify, trim, or weight using \(e(X_i)\).
  \end{itemize}
  \begin{block}{Intuition}
    Instead of matching on the full profile, match on how likely the unit was to be treated.
  \end{block}
\end{frame}

\begin{frame}{The Propensity-Score Theorem}
  \hypertarget{main:pscore-theorem}{}
  If treatment is conditionally independent of potential outcomes given \(X_i\),
  \[
    (Y_{1i},Y_{0i})\indep D_i \mid X_i,
  \]
  then treatment is also conditionally independent given the propensity score:
  \[
    (Y_{1i},Y_{0i})\indep D_i \mid e(X_i).
  \]
  \begin{itemize}
    \item The propensity score is a balancing score.
    \item We can condition on a scalar instead of the whole vector \(X_i\).
    \item This does not mean the propensity-score model can be sloppy.
  \end{itemize}
  \derivbutton{pscore-proof}
\end{frame}

\begin{frame}{Balancing Score Intuition}
  If two units have the same \(e(X)\), then their covariates may differ, but their predicted treatment chance is the same.
  \vspace{0.3em}

  Within a narrow propensity-score bin:
  \[
    \Prob(D_i=1\mid e(X_i)=p)=p.
  \]
  \begin{itemize}
    \item Treatment should not be systematically related to the covariates used to estimate \(e(X)\).
    \item We check this by comparing covariate balance within score strata or after weighting.
    \item If balance is poor, revise the score model or change the design.
  \end{itemize}
  \begin{block}{Beginner trap}
    A high predictive AUC is not the goal. Balance and overlap are the goals.
  \end{block}
\end{frame}

\begin{frame}{Estimating the Propensity Score}
  Common first step:
  \[
    \widehat{e}(X_i)
    =
    \Lambda(X_i'\widehat{\gamma}),
    \qquad
    \Lambda(z)=\frac{\exp(z)}{1+\exp(z)}.
  \]
  \begin{itemize}
    \item Use only pre-treatment variables.
    \item Include variables that affect both treatment and outcome.
    \item Flexible terms can help balance: interactions, squares, splines, or machine learning.
    \item Do not include instruments merely because they predict treatment; they can worsen overlap.
  \end{itemize}
  \begin{block}{The score model is a design model}
    Its purpose is to create comparable groups, not to explain treatment choice beautifully.
  \end{block}
\end{frame}

\begin{frame}{Overlap Diagnostics}
  \hypertarget{main:diagnostics}{}
  After estimating \(\widehat{e}(X)\), inspect:
  \begin{itemize}
    \item Histograms or density plots by treatment status.
    \item Minimum and maximum score in treated and untreated groups.
    \item Fraction of observations near 0 or 1.
    \item Effective sample size after weighting.
    \item Standardized mean differences before and after adjustment.
  \end{itemize}
  \begin{block}{Standardized mean difference}
    \[
      \mathrm{SMD}(X_j)
      =
      \frac{\bar{X}_{j,1}-\bar{X}_{j,0}}{s_j}.
    \]
    Values close to zero indicate better balance.
  \end{block}
\end{frame}

\begin{frame}{Trimming}
  \hypertarget{main:trimming}{}
  Trimming removes observations with weak overlap.
  \[
    a \leq \widehat{e}(X_i) \leq b.
  \]
  Common examples:
  \[
    0.05\leq \widehat{e}(X_i)\leq 0.95
    \qquad \text{or} \qquad
    \max\{\min \widehat{e}_1,\min \widehat{e}_0\}
    \leq
    \widehat{e}_i
    \leq
    \min\{\max \widehat{e}_1,\max \widehat{e}_0\}.
  \]
  \begin{itemize}
    \item Trimming can reduce extrapolation and extreme weights.
    \item It changes the target population.
    \item Always report how many observations were dropped.
  \end{itemize}
  \derivbutton{overlap-trimming}
\end{frame}

\begin{frame}{Propensity-Score Matching}
  Use the scalar score instead of the full covariate vector.
  \[
    d(i,\ell)=|\widehat{e}(X_i)-\widehat{e}(X_\ell)|.
  \]
  \begin{itemize}
    \item Treated units are matched to controls with similar estimated treatment probability.
    \item Calipers are easy: only accept matches within a score distance.
    \item Balance should still be checked on the original covariates, not only on the score.
  \end{itemize}
  \begin{block}{Key limitation}
    Two units with the same score can have different covariate profiles. The theorem is about distributions, not about perfect twins.
  \end{block}
\end{frame}

\section{Inverse Probability Weighting}

\begin{frame}{IPW: Reweighting the Sample}
  \hypertarget{main:ipw}{}
  Inverse probability weighting gives more weight to observations that are rare for their covariates.
  \[
    w_i^{ATE}
    =
    \frac{D_i}{e(X_i)}
    +
    \frac{1-D_i}{1-e(X_i)}.
  \]
  \begin{itemize}
    \item Treated unit with low \(e(X)\): rare treated type, large weight.
    \item Control unit with high \(e(X)\): rare control type, large weight.
    \item Extreme scores create extreme weights.
  \end{itemize}
  \begin{block}{Intuition}
    Weighting reconstructs a pseudo-population where treatment no longer depends on observed covariates.
  \end{block}
  \derivbutton{ipw-proof}
\end{frame}

\begin{frame}{IPW Estimator for the ATE}
  A normalized, or Hajek-style, version is
  \[
    \widehat{\ATE}_{IPW}
    =
    \frac{\sum_i D_iY_i/\widehat{e}(X_i)}
         {\sum_i D_i/\widehat{e}(X_i)}
    -
    \frac{\sum_i (1-D_i)Y_i/[1-\widehat{e}(X_i)]}
         {\sum_i (1-D_i)/[1-\widehat{e}(X_i)]}.
  \]
  \begin{itemize}
    \item The treated group is weighted to represent everyone.
    \item The untreated group is weighted to represent everyone.
    \item The difference estimates the population average treatment effect under CIA and overlap.
  \end{itemize}
\end{frame}

\begin{frame}{IPW Estimator for the ATT}
  For the treated population, treated units already represent treated units.
  Controls must be reweighted to look like the treated:
  \[
    w_i^{ATT}
    =
    \begin{cases}
      1, & D_i=1,\\
      \widehat{e}(X_i)/[1-\widehat{e}(X_i)], & D_i=0.
    \end{cases}
  \]
  Then
  \[
    \widehat{\ATT}_{IPW}
    =
    \frac{\sum_i D_iY_i}{\sum_iD_i}
    -
    \frac{\sum_i (1-D_i)w_i^{ATT}Y_i}{\sum_i(1-D_i)w_i^{ATT}}.
  \]
  \begin{itemize}
    \item Controls who looked likely to be treated receive larger weight.
    \item This targets the treated group, not the full population.
  \end{itemize}
\end{frame}

\begin{frame}{Weights Need Diagnostics}
  \small
  \begin{itemize}
    \item Report the largest weights and percentiles of weights.
    \item Check effective sample size:
      \[
        N_{\mathrm{eff}}=\frac{(\sum_i w_i)^2}{\sum_i w_i^2}.
      \]
    \item Compare covariate balance before and after weighting.
    \item Consider trimming or overlap weights when extreme weights dominate.
  \end{itemize}
  \begin{block}{Interpretation problem}
    If 2,000 observations behave like 60 effective observations after weighting, the design is fragile even if the point estimate looks precise.
  \end{block}
\end{frame}

\section{Practical Workflow}

\begin{frame}{A Practical Workflow}
  \begin{enumerate}
    \item Define the target estimand: ATE, ATT, or an overlap population effect.
    \item Choose pre-treatment covariates using economic reasoning.
    \item Inspect raw treatment-control differences in covariates.
    \item Estimate a matching or propensity-score design.
    \item Check balance, overlap, dropped units, and weights.
    \item Estimate treatment effects.
    \item Report uncertainty and sensitivity to design choices.
  \end{enumerate}
  \begin{block}{Order matters}
    Diagnostics should come before celebrating the treatment-effect estimate.
  \end{block}
\end{frame}

\begin{frame}{Python Example for This Lecture}
  The companion script is:
  \[
    \texttt{lectures/code/lecture-6-matching-propensity.py}.
  \]
  It simulates a setting where:
  \begin{itemize}
    \item treatment is more likely for some observed covariate profiles;
    \item treatment effects are heterogeneous;
    \item the true ATE and ATT are known from simulated potential outcomes;
    \item naive, regression-adjusted, nearest-neighbor matched, and propensity-score weighted estimates are compared.
  \end{itemize}
  \begin{block}{Purpose}
    The script is not a recipe for publication. It is a controlled lab for seeing selection bias, matching, overlap, and weighting.
  \end{block}
\end{frame}

\begin{frame}{What the Script Demonstrates}
  \small
  \begin{tabular}{>{\raggedright\arraybackslash}p{0.28\linewidth}>{\raggedright\arraybackslash}p{0.55\linewidth}}
    \toprule
    Estimate & What students should look for \\
    \midrule
    Naive treated-control gap & Combines treatment effect and selection bias. \\
    Regression adjusted & Works well when the conditional mean is modeled adequately. \\
    Nearest-neighbor matching & Targets the ATT by imputing missing untreated outcomes for treated units. \\
    Propensity-score IPW & Reweights observations to target the ATE or ATT. \\
    Trimmed IPW & Shows how overlap restrictions can stabilize weights and change the sample. \\
    \bottomrule
  \end{tabular}
  \vspace{0.6em}

  The same data can give different estimates because methods target different populations and use different modeling assumptions.
\end{frame}

\begin{frame}{How to Present a Matching or Propensity-Score Result}
  Minimum reporting:
  \begin{itemize}
    \item Causal estimand and identifying assumption.
    \item Covariates used and why they are pre-treatment.
    \item Matching method or propensity-score model.
    \item Balance before and after adjustment.
    \item Overlap plots or score ranges.
    \item Number of observations dropped or reused.
    \item Main estimate, standard error, and sensitivity checks.
  \end{itemize}
  \begin{block}{Reader question}
    Would I believe treated and comparison units are exchangeable after this adjustment?
  \end{block}
\end{frame}

\begin{frame}{Common Failure Modes}
  \small
  \begin{itemize}
    \item Matching on variables affected by treatment.
    \item Treating a propensity-score regression table as evidence of causal validity.
    \item Reporting no balance table or overlap diagnostic.
    \item Estimating an ATE when there is overlap only for the treated population.
    \item Keeping extreme weights because the point estimate is convenient.
    \item Using matching to hide a weak research design.
  \end{itemize}
  \begin{block}{Bottom line}
    Matching and propensity scores can make selection-on-observables designs more transparent, but they do not create unobserved counterfactuals.
  \end{block}
\end{frame}

\begin{frame}{Takeaways}
  \begin{enumerate}
    \item Matching compares treated and untreated units with similar pre-treatment covariates.
    \item Exact matching is clear but usually infeasible in high dimensions.
    \item Nearest-neighbor matching is imputation of missing potential outcomes.
    \item Propensity scores reduce covariates to a scalar balancing score.
    \item Overlap is empirical and must be diagnosed.
    \item IPW creates a pseudo-population but can be unstable with extreme scores.
    \item The causal assumption remains selection on observables.
  \end{enumerate}
\end{frame}

\begin{frame}{Sources and Further Reading}
  \begin{itemize}
    \item Matheus Facure Alves, \emph{Causal Inference for the Brave and True}, Chapter 10, Matching:
      \url{https://matheusfacure.github.io/python-causality-handbook/10-Matching.html}
    \item Matheus Facure Alves, \emph{Causal Inference for the Brave and True}, Chapter 11, Propensity Score:
      \url{https://matheusfacure.github.io/python-causality-handbook/11-Propensity-Score.html}
    \item Rosenbaum, P. R. and Rubin, D. B. (1983). ``The central role of the propensity score in observational studies for causal effects.''
    \item Angrist, J. D. and Pischke, J.-S. (2009). \emph{Mostly Harmless Econometrics}.
  \end{itemize}
\end{frame}

\appendix

\begin{frame}{Appendix Map}
  \begin{itemize}
    \item A1: Conditional independence and matching identification.
    \item A2: Exact matching as the g-formula.
    \item A3: Nearest-neighbor matching as imputation.
    \item A4: Bias from imperfect matches.
    \item A5: Propensity-score theorem.
    \item A6: IPW identities.
    \item A7: Overlap, trimming, and target populations.
  \end{itemize}
\end{frame}

\begin{frame}{Appendix A1: CIA and Matching Identification}
  \hypertarget{app:cia-matching}{}
  Under conditional independence,
  \[
    (Y_{1i},Y_{0i})\indep D_i \mid X_i.
  \]
  Therefore, for \(d\in\{0,1\}\),
  \[
    \E[Y_{di}\mid X_i=x,D_i=1]
    =
    \E[Y_{di}\mid X_i=x,D_i=0]
    =
    \E[Y_{di}\mid X_i=x].
  \]
  Since observed outcomes reveal the realized potential outcome,
  \[
    \E[Y_i\mid X_i=x,D_i=d]
    =
    \E[Y_{di}\mid X_i=x,D_i=d]
    =
    \E[Y_{di}\mid X_i=x].
  \]
  \backbutton{cia}
\end{frame}

\begin{frame}{Appendix A2: Exact Matching as the G-Formula}
  \hypertarget{app:exact-matching}{}
  For the ATT,
  \[
    \ATT
    =
    \E[Y_{1i}-Y_{0i}\mid D_i=1].
  \]
  Apply iterated expectations over \(X_i\):
  \[
    \ATT
    =
    \sum_x
    \left\{
      \E[Y_{1i}\mid X_i=x,D_i=1]
      -
      \E[Y_{0i}\mid X_i=x,D_i=1]
    \right\}
    \Prob(X_i=x\mid D_i=1).
  \]
  By CIA and consistency,
  \[
    \ATT
    =
    \sum_x
    \left\{
      \E[Y_i\mid X_i=x,D_i=1]
      -
      \E[Y_i\mid X_i=x,D_i=0]
    \right\}
    \Prob(X_i=x\mid D_i=1).
  \]
  \backbutton{setup}
\end{frame}

\begin{frame}{Appendix A3: Nearest-Neighbor Matching as Imputation}
  \hypertarget{app:nn-imputation}{}
  For each treated unit \(i\), define weights over controls:
  \[
    W_{i\ell}
    =
    \begin{cases}
      1/K, & \ell\in\mathcal{N}_K(i;D=0),\\
      0, & \text{otherwise}.
    \end{cases}
  \]
  The imputed untreated outcome is
  \[
    \widehat{Y}_{0i}=\sum_{\ell:D_\ell=0}W_{i\ell}Y_\ell.
  \]
  The ATT estimator is
  \[
    \widehat{\ATT}_{M}
    =
    \frac{1}{N_1}\sum_{i:D_i=1}
    \left(
      Y_i-\sum_{\ell:D_\ell=0}W_{i\ell}Y_\ell
    \right).
  \]
  \backbutton{imputation}
\end{frame}

\begin{frame}{Appendix A4: Bias from Imperfect Matches}
  \hypertarget{app:matching-bias}{}
  Let \(m_0(x)=\E[Y_{0i}\mid X_i=x]\). For a treated unit \(i\), its matched control has covariates \(X_{m(i)}\).
  The expected imputation error is approximately
  \[
    \E[Y_{0,m(i)}-Y_{0i}\mid X_i,X_{m(i)}]
    =
    m_0(X_{m(i)})-m_0(X_i).
  \]
  If \(m_0(\cdot)\) is smooth,
  \[
    m_0(X_{m(i)})-m_0(X_i)
    \approx
    \nabla m_0(X_i)'(X_{m(i)}-X_i).
  \]
  \begin{itemize}
    \item Bias shrinks when matches are close.
    \item In high dimensions, nearest neighbors are often not close.
  \end{itemize}
  \backbutton{biasvariance}
\end{frame}

\begin{frame}{Appendix A5: Proof of the Propensity-Score Theorem}
  \hypertarget{app:pscore-proof}{}
  Let \(e(X_i)=\Prob(D_i=1\mid X_i)\). First, \(D_i\indep X_i\mid e(X_i)\) in the balancing-score sense because
  \[
    \Prob(D_i=1\mid X_i,e(X_i))=e(X_i)
  \]
  and
  \[
    \Prob(D_i=1\mid e(X_i))=e(X_i).
  \]
  If \((Y_{1i},Y_{0i})\indep D_i\mid X_i\), then for \(d\in\{0,1\}\),
  \[
    \Prob(D_i=1\mid Y_{di},X_i)=\Prob(D_i=1\mid X_i)=e(X_i).
  \]
  Averaging over all \(X_i\) values with the same score gives
  \[
    \Prob(D_i=1\mid Y_{di},e(X_i))=e(X_i)
    =
    \Prob(D_i=1\mid e(X_i)).
  \]
  Hence \(Y_{di}\indep D_i\mid e(X_i)\).
  \backbutton{pscore-theorem}
\end{frame}

\begin{frame}{Appendix A6: IPW Identities}
  \hypertarget{app:ipw-proof}{}
  Under CIA and overlap,
  \[
    \E\left[\frac{D_iY_i}{e(X_i)}\right]
    =
    \E\left[
      \frac{D_iY_{1i}}{e(X_i)}
    \right].
  \]
  Condition on \(X_i\):
  \[
    \E\left[
      \frac{D_iY_{1i}}{e(X_i)}
      \mid X_i
    \right]
    =
    \frac{\E[D_i\mid X_i]\E[Y_{1i}\mid X_i]}{e(X_i)}
    =
    \E[Y_{1i}\mid X_i].
  \]
  Therefore
  \[
    \E\left[\frac{D_iY_i}{e(X_i)}\right]=\E[Y_{1i}].
  \]
  Similarly,
  \[
    \E\left[\frac{(1-D_i)Y_i}{1-e(X_i)}\right]=\E[Y_{0i}].
  \]
  \backbutton{ipw}
\end{frame}

\begin{frame}{Appendix A7: Overlap, Trimming, and Target Populations}
  \hypertarget{app:overlap-trimming}{}
  Suppose we keep only observations with \(a\leq e(X_i)\leq b\). The estimand becomes
  \[
    \E[Y_{1i}-Y_{0i}\mid a\leq e(X_i)\leq b].
  \]
  This may be more credible than the full ATE when overlap is weak.
  \vspace{0.4em}

  For a weighted estimator, extreme scores inflate variance:
  \[
    \Var\left(\frac{D_iY_i}{e(X_i)}\right)
  \]
  becomes large when \(e(X_i)\) is close to zero.
  \begin{block}{Conclusion}
    Trimming is not just a numerical fix. It defines the population for which the data contain credible treated-control comparisons.
  \end{block}
  \backbutton{overlap}
\end{frame}

\end{document}
